%% cvforge-template v1
%% sections: experience, education, skills
%% supports: person.phone, person.links, experience.location, skills.groups
%% max-bullets-per-role: 5
%% engine: lualatex
%
% Awesome-CV resume.tex (LPPL 1.3c, Byungjin Park), converted to the cvforge contract. Class,
% fonts (Source Sans 3, Roboto, Font Awesome 6), colours and layout are the original's.
% Sections without profile data (summary, honors, certificates) are removed; skills use the
% class's own cvskills environment. Needs tests/fixtures/templates/provision.sh for its fonts.
%!TEX TS-program = lualatex
%!TEX encoding = UTF-8 Unicode
% Awesome CV LaTeX Template for CV/Resume
%
% This template has been downloaded from:
% https://github.com/posquit0/Awesome-CV
%
% Author:
% Claud D. Park <posquit0.bj@gmail.com>
% http://www.posquit0.com
%
% Template license:
% CC BY-SA 4.0 (https://creativecommons.org/licenses/by-sa/4.0/)
%


%-------------------------------------------------------------------------------
% CONFIGURATIONS
% Build from the repository root with: latexmk examples/resume.tex
%-------------------------------------------------------------------------------
% A4 paper size by default, use 'letterpaper' for US letter
\documentclass[11pt, a4paper]{awesome-cv}

% Configure page margins with geometry
\geometry{left=1.4cm, top=.8cm, right=1.4cm, bottom=1.8cm, footskip=.5cm}

% Color for highlights
% Awesome Colors: awesome-emerald, awesome-skyblue, awesome-red, awesome-pink, awesome-orange
%                 awesome-nephritis, awesome-concrete, awesome-darknight
\colorlet{awesome}{awesome-red}
% Uncomment if you would like to specify your own color
% \definecolor{awesome}{HTML}{3E6D9C}

% Colors for text
% Uncomment if you would like to specify your own color
% \definecolor{darktext}{HTML}{414141}
% \definecolor{text}{HTML}{333333}
% \definecolor{graytext}{HTML}{5D5D5D}
% \definecolor{lighttext}{HTML}{999999}
% \definecolor{sectiondivider}{HTML}{5D5D5D}

% Set false if you don't want to highlight section with awesome color
\setbool{acvSectionColorHighlight}{true}

% If you would like to change the social information separator from a pipe (|) to something else
\renewcommand{\acvHeaderSocialSep}{\quad\textbar\quad}


%-------------------------------------------------------------------------------
%	PERSONAL INFORMATION
%	Comment any of the lines below if they are not required
%-------------------------------------------------------------------------------
% Available options: circle|rectangle,edge/noedge,left/right
%-------------------------------------------------------------------------------
%	PERSONAL INFORMATION (cvforge: filled from the profile)
%-------------------------------------------------------------------------------
\name{Jane Doe}{}
\position{Full-Stack Engineer}
\address{Chiang Mai, TH}
\mobile{+66 80 000 0000}
\email{jane@example.com}
\github{janedoe}
\linkedin{janedoe}
%% cvforge: the sample \quote is removed; the profile has no quote.

\begin{document}

\makecvheader[C]

\makecvfooter
  {\today}
  {Jane Doe~~~·~~~Résumé}
  {\thepage}

%-------------------------------------------------------------------------------
%	EXPERIENCE (as in examples/resume/experience.tex)
%-------------------------------------------------------------------------------
\cvsection{Work Experience}

\begin{cventries}

  \cventry
    {Technical Writer} % Job title
    {Acme Software} % Organization
    {Bangkok, TH} % Location
    {Jun 2025 - Present} % Date(s)
    {
      \begin{cvitems} % Description(s) of tasks/responsibilities
        \item {Built a Python pipeline migrating 400 knowledge-base articles to DocFX.} % id:exp-01.b2
        \item {Cut docs build time 40\% by caching DocFX output in GitHub Actions.} % id:exp-01.b4
        \item {Wrote C\# code samples for every public API \& kept them compiling in CI.} % id:exp-01.b3
        \item {Owned end-to-end documentation for a .NET PDF library used by 2,000 teams.} % id:exp-01.b1
        \item {Typeset the printed API reference in LaTeX from the same Markdown sources.} % id:exp-01.b6
      \end{cvitems}
    }

  \cventry
    {Junior Full-Stack Developer} % Job title
    {Lanna Digital Co., Ltd.} % Organization
    {Chiang Mai, TH} % Location
    {Jan 2024 - May 2025} % Date(s)
    {
      \begin{cvitems} % Description(s) of tasks/responsibilities
        \item {Shipped a REST API in Node.js and PostgreSQL serving 1M requests a day.} % id:exp-02.b1
        \item {Built a React and TypeScript dashboard used by 30 field staff for daily reporting.} % id:exp-02.b2
        \item {Containerised three services with Docker, cutting onboarding from two days to one hour.} % id:exp-02.b3
        \item {Added GitHub Actions pipelines running tests on every pull request.} % id:exp-02.b4
        \item {Wrote the internal API guide that new developers read in their first week.} % id:exp-02.b5
      \end{cvitems}
    }

\end{cventries}

%-------------------------------------------------------------------------------
%	EDUCATION (as in examples/resume/education.tex)
%-------------------------------------------------------------------------------
\cvsection{Education}

\begin{cventries}

  \cventry
    {B.Eng. Computer Innovation Engineering} % Degree
    {KMITL, Bangkok} % Institution
    {} % Location
    {Aug 2020 - May 2024} % Date(s)
    {}

\end{cventries}

%-------------------------------------------------------------------------------
%	SKILLS (the class's cvskills environment)
%-------------------------------------------------------------------------------
\cvsection{Skills}

\begin{cvskills}
  \cvskill{Frontend}{TypeScript, React}
  \cvskill{Backend}{REST APIs, Node.js, PostgreSQL, Python, .NET / C\#}
  \cvskill{Tooling}{GitHub Actions, Technical documentation, Docker, LaTeX}
\end{cvskills}

\end{document}
